\chapter{Replicarea relațiilor}
\label{sec:replicare}

\section*{a. Relații replicate în proiect}

Următoarele relații au fost replicate la nivelul sistemului: \textit{STATUS}, \textit{PRIORITATE}, \textit{CATEGORIE}, \textit{TAG} și \textit{TOPIC}.

\section*{b. Motivație și avantaje}

Aceste relații au un volum redus de date și sunt utilizate frecvent de toate componentele sistemului distribuit. Replicarea lor oferă avantaje semnificative, precum:

\begin{itemize}
    \item \textbf{Acces local rapid la date.}
    \item \textbf{Îmbunătățirea performanței interogărilor} (elimină necesitatea accesului prin rețea pentru obținerea detaliilor de bază).
    \item \textbf{Creșterea disponibilității sistemului.}
\end{itemize}

\section*{c. Costuri de sincronizare}

De asemenea, costul sincronizării între nodurile sistemului distribuit este redus, deoarece datele din aceste relații (fiind în general de tip nomenclator) sunt modificate foarte rar.